% GENERATED FILE. Edit canonical lessons, then run npm run generate:books.
% canonical-source-hash: fnv1a64:62c707d12af359f3
% canonical-lessons: LA-C101-surgo, LA-C101-erro, LA-C101-fugio, LA-C101-statim, LA-C101-deinde

\chapter{Get Up, Wander, Run Away, At Once, Then}
\label{ch:la-get-up-wander-run-away-at-once-then}
\hlchaptermodality{101}

\begin{chapteropening}
\textbf{By the end of this chapter:} \emph{I can say to rise, to wander and to flee, and at once and then.}

Complete the last lesson of chapter 101: I can say to rise, to wander and to flee, and at once and then.
\end{chapteropening}

\section[surgō, surgere]{surgō, surgere — to rise, to get up}

\label{lesson:LA-C101-surgo}



\begin{quote}
Before the new one: say the Latin for to touch, then the Latin for to pull.
\end{quote}

\subsection*{You'll want to know: surgō, surgere}
\textbf{surgō, surgere} — \textquotedblleft{}to rise, to get up\textquotedblright{}. English \emph{surge}.

\begin{cousinweb}
\textbf{surgō} is contracted from \emph{surrigō}, \textquotedblleft{}to raise up\textquotedblright{}, built from \emph{sub}, \textquotedblleft{}up from under\textquotedblright{}, and \emph{regō}, \textquotedblleft{}to direct, to keep straight\textquotedblright{} (Lewis \& Short).
\end{cousinweb}

\begin{grammarlens}[title={What you've built}]
One more word for this chapter.
\end{grammarlens}

\subsection*{Your turn}
\begin{itemize}
\raggedright
  \item \emph{Say it:} \emph{surgō}
  \item \emph{Say it:} \emph{surgō}, once more
  \item \emph{Say it:} say \emph{surgō}
  \item \emph{Recall:} say the Latin for to touch, then the Latin for to pull, then say \emph{surgō} again
\end{itemize}

\begin{tcolorbox}[breakable,colback=teal!4,colframe=teal!35!black,title={Before you move on}]
What does surgō, surgere mean? (\textquotedblleft{}To rise, to get up\textquotedblright{}.) Say it once more.
\end{tcolorbox}
\section[errō, errāre]{errō, errāre — to wander; to be wrong}

\label{lesson:LA-C101-erro}



\begin{quote}
Before the new one: say the Latin for to pull, then the Latin for to rise.
\end{quote}

\subsection*{You'll want to know: errō, errāre}
\textbf{errō, errāre} — \textquotedblleft{}to wander; to be wrong\textquotedblright{}. English \emph{err} and \emph{error}.

\begin{cousinweb}
\textbf{errō} is inherited, usually traced to Proto-Indo-European \emph{*h1ers-}, \textquotedblleft{}to be in motion, to wander\textquotedblright{} (EtymDB 2.1); German \emph{irren}, \textquotedblleft{}to err\textquotedblright{}, is a cousin (Lewis \& Short).
\end{cousinweb}

\begin{grammarlens}[title={What you've built}]
One more word for this chapter.
\end{grammarlens}

\subsection*{Your turn}
\begin{itemize}
\raggedright
  \item \emph{Say it:} \emph{errō}
  \item \emph{Say it:} \emph{errō}, once more
  \item \emph{Say it:} say \emph{errō}
  \item \emph{Recall:} say the Latin for to pull, then the Latin for to rise, then say \emph{errō} again
\end{itemize}

\begin{tcolorbox}[breakable,colback=teal!4,colframe=teal!35!black,title={Before you move on}]
What does errō, errāre mean? (\textquotedblleft{}To wander; to be wrong\textquotedblright{}.) Say it once more.
\end{tcolorbox}
\section[fugiō, fugere]{fugiō, fugere — to flee, to run away}

\label{lesson:LA-C101-fugio}



\begin{quote}
Before the new one: say the Latin for to rise, then the Latin for to wander; to be wrong.
\end{quote}

\subsection*{You'll want to know: fugiō, fugere}
\textbf{fugiō, fugere} — \textquotedblleft{}to flee, to run away\textquotedblright{}. English \emph{fugitive}.

\begin{cousinweb}
\textbf{fugiō} is inherited, usually traced to Proto-Indo-European \emph{*bhewg-}, \textquotedblleft{}to flee\textquotedblright{} (EtymDB 2.1); Greek \emph{pheugō}, \textquotedblleft{}to flee\textquotedblright{}, is a cousin (Lewis \& Short).
\end{cousinweb}

\begin{grammarlens}[title={What you've built}]
One more word for this chapter.
\end{grammarlens}

\subsection*{Your turn}
\begin{itemize}
\raggedright
  \item \emph{Say it:} \emph{fugiō}
  \item \emph{Say it:} \emph{fugiō}, once more
  \item \emph{Say it:} say \emph{fugiō}
  \item \emph{Recall:} say the Latin for to rise, then the Latin for to wander; to be wrong, then say \emph{fugiō} again
\end{itemize}

\begin{tcolorbox}[breakable,colback=teal!4,colframe=teal!35!black,title={Before you move on}]
What does fugiō, fugere mean? (\textquotedblleft{}To flee, to run away\textquotedblright{}.) Say it once more.
\end{tcolorbox}
\section[statim]{statim — at once, immediately}

\label{lesson:LA-C101-statim}



\begin{quote}
Before the new one: say the Latin for to wander; to be wrong, then the Latin for to flee.
\end{quote}

\subsection*{You'll want to know: statim}
\textbf{statim} — \textquotedblleft{}at once, immediately\textquotedblright{}.

\begin{cousinweb}
\textbf{statim} is built on the root of \emph{stō}, \textquotedblleft{}to stand\textquotedblright{} (Lewis \& Short): first \textquotedblleft{}standing firm, on the spot\textquotedblright{}, then \textquotedblleft{}at once\textquotedblright{}.
\end{cousinweb}

\begin{grammarlens}[title={What you've built}]
One more word for this chapter.
\end{grammarlens}

\subsection*{Your turn}
\begin{itemize}
\raggedright
  \item \emph{Say it:} \emph{statim}
  \item \emph{Say it:} \emph{statim}, once more
  \item \emph{Say it:} say \emph{statim}
  \item \emph{Recall:} say the Latin for to wander; to be wrong, then the Latin for to flee, then say \emph{statim} again
\end{itemize}

\begin{tcolorbox}[breakable,colback=teal!4,colframe=teal!35!black,title={Before you move on}]
What does statim mean? (\textquotedblleft{}At once, immediately\textquotedblright{}.) Say it once more.
\end{tcolorbox}
\section[deinde]{deinde — then, next}

\label{lesson:LA-C101-deinde}



\begin{quote}
Before the new one: say the Latin for to flee, then the Latin for at once.
\end{quote}

\subsection*{You'll want to know: deinde}
\textbf{deinde} — \textquotedblleft{}then, next\textquotedblright{}.

\begin{cousinweb}
\textbf{deinde} is built from \emph{dē}, \textquotedblleft{}from\textquotedblright{}, and \emph{inde}, \textquotedblleft{}from there, thence\textquotedblright{} (Lewis \& Short): from that point on.
\end{cousinweb}

\begin{grammarlens}[title={What you've built}]
One more word for this chapter.
\end{grammarlens}

\subsection*{Your turn}
\begin{itemize}
\raggedright
  \item \emph{Say it:} \emph{deinde}
  \item \emph{Say it:} \emph{deinde}, once more
  \item \emph{Say it:} say \emph{deinde}
  \item \emph{Recall:} say the Latin for to flee, then the Latin for at once, then say \emph{deinde} again
\end{itemize}

\begin{tcolorbox}[breakable,colback=teal!4,colframe=teal!35!black,title={Before you move on}]
What does deinde mean? (\textquotedblleft{}Then, next\textquotedblright{}.) Say it once more.
\end{tcolorbox}
